% Figure 2. The same hour of work in four records, and what one operation costs
% in each. Same session as Figure 1, same colours, same glyphs.
\newcommand{\FigTwoCaption}{One hour of work with a coding agent, recorded four
ways, with the same request removed from each. The unit a system records decides
which sets a developer can name. Git and Darcs both bottom out at a line, so
removing the caching means selecting lines. Jujutsu and Mercurial add a durable
identity above the line, and that identity belongs to a commit, so a squashed
session carries one identity for all three requests. \sgt{} records one edit per
function together with the request it was made for. The developer supplied that
request and never wrote it down.}

\newcommand{\FigTwoDescription}{A four row comparison. Each row is a card named
after a version control system, drawing the same one hour agent session in that
system's units, with the unit of record stated on the left and the cost of
removing the caching request on the right. The git row draws two commit boxes
filled with coloured diff stripes of four different colours mixed together. The
Darcs and Pijul row draws nine patch tokens in a line with arcs above joining the
patches that share lines, and one dashed arc below joining two patches that depend
on each other through a function call, marked as not recorded. The Jujutsu and
Mercurial row draws one box of mixed stripes, an arrow, and a second box of the
same mixed stripes, labelled with the same change identifier and two different
commit hashes. The sgt row draws four labelled horizontal lines, one per request,
each carrying a mark for every edit filed under that request.}

\begin{tikzpicture}[x=1mm, y=1mm, line cap=round, font=\sffamily]

% ------------------------------- headers ---------------------------------
\node[hd, anchor=west] at (2,79.5)   {UNIT OF RECORD};
\node[hd, anchor=west] at (30,79.5)  {THE SAME HOUR OF WORK, RECORDED FOUR WAYS};
\node[hd, anchor=west] at (117,79.5) {TAKING THE CACHING BACK OUT};

% Every row is one card. The cost of the operation gets a band of its own on
% the right, so the reader can read that column straight down.
% #1 y, #2 height, #3 tab colour, #4 tab text
\newcommand{\cmprow}[4]{%
  \panel{0}{#1}{176}{#2}{#3}{#4}
  \fill[wGrey, rounded corners=0.8mm] (114.5,{#1+1.2}) rectangle (174.6,{#1+#2-1.2});}
% #1 y, #2 text
\newcommand{\rowunit}[2]{%
  \node[anchor=north west, text width=23mm, align=left, inner sep=0, text=sgtInk!75,
        font=\sffamily\fontsize{5.6}{6.5}\selectfont] at (2.2,#1) {#2};}
\newcommand{\rowcost}[2]{%
  \node[anchor=north west, text width=55.5mm, align=left, inner sep=0, text=sgtInk!85,
        font=\sffamily\fontsize{5.7}{6.7}\selectfont] at (116.5,#1) {#2};}

% =============================== git =====================================
\cmprow{55}{19}{sgtGrey}{GIT}
\rowunit{71.4}{a snapshot of the tree. An operation names a line inside a commit.}
\rowcost{71.4}{Select the 6 caching hunks out of 15, in 3 files, across both
  commits, and reverse them. Then read what is left to confirm that nothing
  still calls \fsym{get\_cached}.}

\draw[card, color=sgtInk!18, fill=white] (30,58.2) rectangle (68,70.9);
\node[anchor=north east, inner sep=0, text=sgtGrey,
      font=\ttfamily\fontsize{5}{5.6}\selectfont] at (67.2,70.5) {c39ab2};
\hunk{32.4}{68.5}{9}{sgtGrey}\hunk{32.4}{67.15}{14}{fCache}
\hunk{32.4}{65.8}{12}{fCache}\hunk{32.4}{64.45}{17}{fRate}
\hunk{32.4}{63.1}{16}{fRate}\hunk{32.4}{61.75}{11}{fRate}
\hunk{32.4}{60.4}{19}{fCache}\hunk{32.4}{59.05}{22}{fCache}

\draw[card, color=sgtInk!18, fill=white] (72,58.2) rectangle (110,70.9);
\node[anchor=north east, inner sep=0, text=sgtGrey,
      font=\ttfamily\fontsize{5}{5.6}\selectfont] at (109.2,70.5) {7f10de};
\hunk{74.4}{68.5}{13}{fRetry}\hunk{74.4}{67.15}{18}{fCache}
\hunk{74.4}{65.8}{21}{fRetry}\hunk{74.4}{64.45}{15}{fRetry}
\hunk{74.4}{63.1}{19}{fRetry}\hunk{74.4}{61.75}{10}{fCache}
\hunk{74.4}{60.4}{16}{sgtGrey}

% ========================== Darcs and Pijul ==============================
\cmprow{37}{15.5}{sgtGrey}{DARCS, PIJUL}
\rowunit{50.1}{a patch of added and removed lines, plus the patches whose lines
  it touched.}
\rowcost{50.1}{Unpull the 6 caching patches. \fsym{db.py::query} loses a
  parameter that a caller in another patch still passes, and no recorded
  dependency says so, because the two patches share no line.}

% #1 x centre, #2 colour
\newcommand{\patchtok}[2]{%
  \draw[card, color=#2!85, fill=#2!30, line width=0.35pt]
    ({#1-3.6},44.4) rectangle ({#1+3.6},47.4);}
\patchtok{34}{fRate}\patchtok{42.5}{fRate}\patchtok{51}{fCache}\patchtok{59.5}{fCache}
\patchtok{68}{fCache}\patchtok{76.5}{sgtGrey}\patchtok{85}{fRetry}\patchtok{93.5}{fRetry}
\patchtok{102}{fCache}

\draw[dep] (59.5,47.7) .. controls (63,50.2) and (98.5,50.2) .. (102,47.7);
\draw[dep] (51,47.7) .. controls (54.5,49.7) and (90,49.7) .. (93.5,47.7);
\draw[dep] (68,47.7) .. controls (70,49.2) and (74.5,49.2) .. (76.5,47.7);
\node[acc=sgtInk!70, anchor=south] at (85,50.2) {patches that share lines};
\draw[color=sgtGrey!85, line width=0.35pt, densely dotted]
  (51,44.1) .. controls (56.0,41.3) and (63.0,41.3) .. (68,44.1);
\node[acc=sgtInk!70, anchor=north] at (59.5,42.1)
  {\fsym{handle} calls \fsym{query}. Not recorded.};

% ======================= Jujutsu and Mercurial ===========================
\cmprow{18}{15.5}{sgtGrey}{JUJUTSU, MERCURIAL}
\rowunit{31.1}{a commit whose identity survives being rewritten.}
\rowcost{31.1}{The sitting is one change with one identifier, so the caching has
  no identifier of its own. Split the commit by selecting lines to give the
  caching one.}

\newcommand{\jjbox}[2]{%
  \draw[card, color=sgtInk!18, fill=white] (#1,20.0) rectangle ({#1+30},30.3);
  \hunk{{#1+2.4}}{28.7}{14}{fRate}\hunk{{#1+2.4}}{27.35}{11}{fCache}
  \hunk{{#1+2.4}}{26.0}{19}{fCache}\hunk{{#1+2.4}}{24.65}{16}{fRetry}
  \hunk{{#1+2.4}}{23.3}{22}{fCache}\hunk{{#1+2.4}}{21.95}{12}{sgtGrey}
  \hunk{{#1+2.4}}{20.6}{17}{fRetry}
  \node[anchor=south west, inner sep=0, text=sgtInk!80,
        font=\ttfamily\fontsize{5.2}{6}\selectfont] at (#1,30.7) {#2};}
\jjbox{30}{kxryzmot / 7f10de}
\jjbox{66}{kxryzmot / a19c4b}
\draw[dep] (60.8,25.2) -- (64.6,25.2);
\node[note, anchor=south, text=sgtGrey] at (62.7,25.9) {rewritten};
\annot{98.9}{28.2}{96.6}{30.2}{20}{sgtInk}
\node[anchor=north west, text width=13mm, align=left, inner sep=0, text=sgtInk!70,
      font=\sffamily\fontsize{5.2}{6}\selectfont] at (99.2,29.0)
  {three requests,\\one identifier};

% ================================= sgt ===================================
\cmprow{0}{17}{sgtInk}{\textsf{sgt}}
\rowunit{14.0}{one edit to one function, filed under the request it was made
  for.}
\node[anchor=north west, inner sep=0, text=sgtInk,
      font=\ttfamily\fontsize{5.9}{7}\selectfont] at (116.5,14.0)
  {\$ sgt revert "price cache"};
\node[anchor=north west, text width=55.5mm, align=left, inner sep=0, text=sgtInk!85,
      font=\sffamily\fontsize{5.7}{6.7}\selectfont] at (116.5,10.2)
  {6 edits, 4 functions, 3 files. The 4 edits recorded as built on top of them
   are named before anything is removed.};

\draw[commitcol] (56,2.2) -- (56,14.0);
\draw[commitcol] (86,2.2) -- (86,14.0);
\node[note, text=sgtGrey, anchor=south] at (56,14.1) {c39ab2};
\node[note, text=sgtGrey, anchor=south] at (86,14.1) {7f10de};

\newcommand{\srail}[3]{%
  \draw[rail=#3] (48,#1) -- (112,#1);
  \node[anchor=east, inner sep=0, text=#3,
        font=\sffamily\fontsize{5.8}{6.6}\selectfont] at (46.6,#1) {#2};}
\srail{12.2}{rate limiting}{fRate}
\opadd{56.5}{12.2}{fRate}\opadd{59.9}{12.2}{fRate}\opadd{63.3}{12.2}{fRate}
\opext{66.7}{12.2}{fRate}\ckpt{69.6}{12.2}{fRate}
\srail{9.0}{price cache}{fCache}
\opadd{56.5}{9.0}{fCache}\opadd{59.9}{9.0}{fCache}\opext{63.3}{9.0}{fCache}
\opext{66.7}{9.0}{fCache}\ckpt{69.6}{9.0}{fCache}
\oprew{86.5}{9.0}{fCache}\opext{89.9}{9.0}{fCache}\ckpt{92.8}{9.0}{fCache}
\srail{5.8}{retry policy}{fRetry}
\opadd{86.5}{5.8}{fRetry}\opadd{89.9}{5.8}{fRetry}\opext{93.3}{5.8}{fRetry}
\ckpt{96.2}{5.8}{fRetry}
\srail{2.6}{row accessors}{sgtGrey}
\whoagent{30.0}{2.6}{sgtGrey}
\opmove{56.5}{2.6}{sgtGrey}\opext{59.9}{2.6}{sgtGrey}\opext{63.3}{2.6}{sgtGrey}
\opext{66.7}{2.6}{sgtGrey}\opext{86.5}{2.6}{sgtGrey}

\end{tikzpicture}
